% =============================================================================
% comparacion_energias_2012.tex — Nota comparativa: las ocho energías de la
% Sección 4 de Gutiérrez et al. (2012), "Inelastic Collisions and
% Hypervelocity Impacts at Nanoscopic Level", vs. los cuatro canales alpha
% de este estudio (borrador_rbef_v2.tex).
% Compilar: pdflatex comparacion_energias_2012.tex (dos pasadas)
% =============================================================================
\documentclass[11pt,a4paper]{article}

\usepackage[utf8]{inputenc}
\usepackage[T1]{fontenc}
\usepackage[spanish,es-nodecimaldot]{babel}
\usepackage[margin=2.3cm]{geometry}
\usepackage{amsmath,amssymb}
\usepackage{booktabs}
\usepackage{enumitem}
\usepackage{xcolor}
\usepackage[colorlinks=true,linkcolor=blue!50!black,citecolor=blue!50!black]{hyperref}

\newcommand{\eps}{\varepsilon}
\newcommand{\kb}{k_B}
\newcommand{\lammps}{\textsc{Lammps}}

\title{\textbf{Las energías del rebote: dos formulaciones}\\
\large Comparación entre la partición de energía de
Gutiérrez \emph{et al.} (2012) y la de este estudio}
\author{Felipe Wasaff}
\date{Julio 2026}

\begin{document}
\maketitle

\begin{abstract}
Gutiérrez \emph{et al.} (2012)\cite{Gutierrez2012}, Sección~4
(``Bouncing of a ball over a surface: atomic level study''), definen ocho
energías para seguir la disipación de un cúmulo LJ que rebota sobre una
superficie. Este estudio parte exactamente de ese mismo esquema, pero lo
reorganiza en cuatro fracciones adimensionales $\alpha_i$ que suman $1$ y
se miden en el primer apogeo. Esta nota escribe ambos conjuntos de
ecuaciones lado a lado, las desarrolla paso a paso, y hace explícito el
mapeo entre ellos, incluyendo dos sutilezas de implementación (el reparto
de la energía colisional $U_c$ y la definición de $K_S$) que no son
evidentes de la sola lectura de ambos textos.
\end{abstract}

\tableofcontents

% =============================================================================
\section{Las ocho energías de Gutiérrez et al. (2012)}
\label{sec:2012}
% =============================================================================

El capítulo trabaja con un sistema cerrado (bola + superficie), integrado
con el algoritmo de Beeman, sin termostato: la energía total se conserva
exactamente y ``la pérdida'' de energía es en realidad una
\emph{transferencia} entre los siguientes ocho términos.

\subsection{Potencial de interacción}

\begin{equation}
V(r) = \begin{cases}
4\eps\left[\left(\dfrac{\sigma}{r}\right)^{12}
- c\left(\dfrac{\sigma}{r}\right)^{6}\right] & r < r_c=2.5\,\sigma\\[4pt]
0 & r\ge r_c,
\end{cases}
\label{eq:2012_lj}
\end{equation}
con $c=1$ para pares del mismo tipo (bola--bola, superficie--superficie) y
$c=0$ para pares bola--superficie (repulsión pura). $\eps/\kb=119.8$~K,
$\sigma=3.40$~\AA{} (argón).

\subsection{Energías cinéticas}

\begin{align}
K_B &= \tfrac12 M_B\|\mathbf V_{\rm CM}^B\|^2
    &&\text{cinética traslacional del CM de la bola,}
    \label{eq:2012_KB}\\
K_S &= \tfrac12 M_S\|\mathbf V_{\rm CM}^S\|^2
    &&\text{cinética traslacional del CM de la superficie,}
    \label{eq:2012_KS}\\
T_B &= \sum_{i\in{\rm bola}} \tfrac12 m_i\|\mathbf v_i-\mathbf V_{\rm CM}^B\|^2
    &&\text{energía térmica (vibracional) de la bola,}
    \label{eq:2012_TB}\\
T_S &= \sum_{i\in{\rm superf.}} \tfrac12 m_i\|\mathbf v_i-\mathbf V_{\rm CM}^S\|^2
    &&\text{energía térmica de la superficie.}
    \label{eq:2012_TS}
\end{align}

\subsection{Energía potencial: cuatro términos}

\begin{equation}
U = V_B + V_S + U_c + U_F,
\label{eq:2012_U}
\end{equation}
\begin{align}
V_B &= \sum_{\text{pares bola}} V(r) &&\text{potencial vibracional de la bola,}
    \label{eq:2012_VB}\\
V_S &= \sum_{\text{pares superficie}} V(r) &&\text{potencial vibracional de la superficie,}
    \label{eq:2012_VS}\\
U_c &= \sum_{\text{pares bola--superficie}} V(r) &&\text{energía colisional,}
    \label{eq:2012_Uc}\\
U_F &= M_B\,a\,z_{\rm CM} &&\text{potencial del campo de fuerza (gravedad),}
    \label{eq:2012_UF}
\end{align}
con $a=0.026\,a_0$ la aceleración impuesta. Como el sistema es cerrado y
newtoniano,
\[
K_B+K_S+T_B+T_S+V_B+V_S+U_c+U_F = \text{const.}
\]
en todo instante ---no solo en el apogeo---, lo que permite graficar la
evolución temporal completa de cada término (Fig.~9 del capítulo) y usar
el ancho del pico de $U_c(t)$ como definición operacional del
\emph{tiempo de contacto}.

\subsection{Desarrollo explícito de los cuatro términos}
\label{sec:desarrollo_2012}

Las ecs.~(\ref{eq:2012_VB})--(\ref{eq:2012_Uc}) son sumas dobles sobre
pares de átomos; escribirlas con índices explícitos deja ver por qué son
tres piezas disjuntas de la misma suma total. Sea $\mathcal B=\{1,\ldots,
N\}$ el conjunto de átomos de la bola y $\mathcal S$ el de la superficie,
con $r_{ij}=\|\vec r_i-\vec r_j\|$. Sustituyendo el potencial
ec.~(\ref{eq:2012_lj}) (con $c=1$ para pares del mismo tipo,
$c=0$ para pares bola--superficie):
\begin{align}
V_B &= \sum_{\substack{i<j\\ i,j\in\mathcal B}}
    4\eps\left[\left(\frac{\sigma}{r_{ij}}\right)^{12}
    -\left(\frac{\sigma}{r_{ij}}\right)^{6}\right],
    \label{eq:2012_VB_exp}\\
V_S &= \sum_{\substack{i<j\\ i,j\in\mathcal S}}
    4\eps\left[\left(\frac{\sigma}{r_{ij}}\right)^{12}
    -\left(\frac{\sigma}{r_{ij}}\right)^{6}\right],
    \label{eq:2012_VS_exp}\\
U_c &= \sum_{i\in\mathcal B}\sum_{j\in\mathcal S}
    4\eps\left(\frac{\sigma}{r_{ij}}\right)^{12},
    \label{eq:2012_Uc_exp}
\end{align}
todas restringidas a $r_{ij}<r_c=2.5\,\sigma$ (fuera de ese radio, cada
término es cero). Nótese que $U_c$, ec.~(\ref{eq:2012_Uc_exp}), no lleva
el término atractivo $-(\sigma/r_{ij})^6$: el par bola--superficie usa
$c=0$ (repulsión pura), de modo que $U_c\ge0$ siempre, mientras que
$V_B$ y $V_S$ sí pueden ser negativos (el cúmulo y la losa son estables
por la parte atractiva de Lennard-Jones). La suma triple
$V_B+V_S+U_c$ recorre \emph{cada} par de átomos del sistema exactamente
una vez: los pares $\mathcal B$--$\mathcal B$ en $V_B$, los pares
$\mathcal S$--$\mathcal S$ en $V_S$, y los pares cruzados
$\mathcal B$--$\mathcal S$ en $U_c$ --- no hay traslape ni pares
faltantes, lo que garantiza que ec.~(\ref{eq:2012_U}) sea exactamente el
potencial total del sistema y no una aproximación.

El cuarto término, $U_F$, no es una suma sobre pares sino la energía
potencial de la fuerza externa constante $\vec F=-M_Ba\,\hat z$ que actúa
solo sobre la bola (simulando gravedad). Integrando
$U_F(z)=-\int_0^z F_z\,dz'=\int_0^z M_Ba\,dz'=M_Ba\,z$ y evaluando en
$z=z_{\rm CM}$, la coordenada del centro de masa de la bola, se recupera
la ec.~(\ref{eq:2012_UF}). En unidades no reducidas, $a$ hace las veces de
la aceleración de gravedad $g$, así que $U_F=M_Bgz_{\rm CM}$ es,
literalmente, la energía potencial gravitatoria de toda la vida.

% =============================================================================
\section{Los cuatro canales de este estudio}
\label{sec:nuestro}
% =============================================================================

Este estudio no grafica energías en función del tiempo continuo, sino que
las evalúa en un único instante ---el primer apogeo, cuando la bola vuelve
a alcanzar velocidad vertical nula tras separarse--- y las normaliza por
la energía de impacto $E_0=Mg\,h_0^{\rm ef}$, de modo que los canales sean
fracciones adimensionales comparables entre corridas de $N$ y $h_0$
distintos:
\begin{align}
    \alpha_{\rm cm}     &= \frac{h_1}{h_0^{\rm ef}} = e_1^2
        &&\text{energía recuperada por el CM,}
        \label{eq:acm}\\
    \alpha_{T_B}        &= \frac{T_B^{(1)}}{E_0}
        &&\text{KE vibracional interna de la bola,}
        \label{eq:atb}\\
    \alpha_{\rm surf}   &= \frac{\Delta(K_S + V_{\rm surf})}{E_0}
        &&\text{fonones de la superficie,}
        \label{eq:asurf}\\
    \alpha_{\Delta V_B} &= \frac{V_B^{(1)}-V_B^{(0)}}{E_0}
        &&\text{calentamiento potencial de la bola,}
        \label{eq:avb}
\end{align}
donde el superíndice $(1)$ indica ``evaluado en el primer apogeo'' y
$(0)$, ``en el instante inicial''. La conservación de energía garantiza
\begin{equation}
\alpha_{\rm cm}+\alpha_{T_B}+\alpha_{\rm surf}+\alpha_{\Delta V_B} = 1
\label{eq:cierre}
\end{equation}
salvo un quinto término pequeño (la KE lateral del CM cuando el rebote no
es perfectamente vertical), que en las tablas se absorbe dentro de
$\alpha_{\Delta V_B}$ reportándola \emph{por cierre},
$\alpha_{\Delta V_B}=1-\alpha_{\rm cm}-\alpha_{T_B}-\alpha_{\rm surf}$, en
vez de calcularla directamente de $V_B$.

% =============================================================================
\section{Mapeo entre ambas formulaciones}
\label{sec:mapeo}
% =============================================================================

\begin{center}
\begin{tabular}{lll}
\toprule
Gutiérrez 2012 & Este estudio & Relación \\
\midrule
$V(r)$, ec.~(\ref{eq:2012_lj}) & $U_{\rm LJ}(r)$ & idéntica; $c\to c_{ij}$ \\
$T_B$, ec.~(\ref{eq:2012_TB}) & $T_B^{(1)}$ en $\alpha_{T_B}$ & idéntica, evaluada en el apogeo \\
$V_B$, ec.~(\ref{eq:2012_VB}) & $V_B^{(1)}-V_B^{(0)}$ en $\alpha_{\Delta V_B}$ & idéntica; nosotros restamos el valor inicial \\
$K_B$, ec.~(\ref{eq:2012_KB}) & $\alpha_{\rm cm}=h_1/h_0^{\rm ef}$ & equivalente: $K_B^{(1)}=\alpha_{\rm cm}E_0$ \\
$U_c$, ec.~(\ref{eq:2012_Uc}) & --- (no aparece en el apogeo) & $U_c\to0$ una vez separada la bola \\
$U_F$, ec.~(\ref{eq:2012_UF}) & absorbido en la definición de $E_0=Mg h_0^{\rm ef}$ & mismo término, contabilizado distinto \\
$K_S+T_S$ & $\alpha_{\rm surf}$ (vía \texttt{KS}) & ver nota (a) abajo \\
$V_S$ & $\alpha_{\rm surf}$ (vía \texttt{PE\_surf}) & ver nota (b) abajo \\
\bottomrule
\end{tabular}
\end{center}

\textbf{Nota (a) --- $K_S$ vs.\ $K_S+T_S$.} El capítulo separa
formalmente la traslación del centro de masa de la superficie ($K_S$,
ec.~\ref{eq:2012_KS}) de su energía térmica interna ($T_S$,
ec.~\ref{eq:2012_TS}). Nuestro código \texttt{in.bounce\_cluster} calcula
en cambio la energía cinética \emph{total} del grupo superficie
(\texttt{compute surf\_ke}, sin restar el CM), es decir mide
$K_S^{\rm paper}+T_S^{\rm paper}$ de una sola vez. Esto no introduce error
en nuestro caso porque las capas inferiores de la losa están fijas
(\texttt{setforce~0~0~0}): la superficie no adquiere momento neto
apreciable, así que $K_S^{\rm paper}\approx0$ y nuestra variable
\texttt{KS} es, en la práctica, $T_S^{\rm paper}$.

\textbf{Nota (b) --- el reparto de $U_c$ entre $V_B$ y $V_S$.} El cómputo
\texttt{pe/atom} de \lammps{} reparte la energía de cada par
\emph{50/50} entre sus dos átomos. Sumar esa cantidad sobre todos los
átomos de la bola da entonces $V_B^{\rm paper}+\tfrac12 U_c$, no
$V_B^{\rm paper}$ puro; por eso nuestro código corrige explícitamente
\verb|VB = c_pe_ball - 0.5*c_Uc|, restando la mitad de $U_c$ que le
corresponde a la bola. La misma suma sobre la superficie
(\texttt{PE\_surf}) \emph{no} se corrige, así que
$\texttt{PE\_surf}=V_S^{\rm paper}+\tfrac12 U_c$. Esto sería un problema
si evaluáramos estas cantidades durante la colisión, pero como
$\alpha_{\rm surf}$ y $\alpha_{\Delta V_B}$ se miden en el primer apogeo
---cuando la bola ya está separada y $U_c\approx0$ por definición
(ec.~\ref{eq:2012_Uc} solo suma pares bola--superficie dentro de
$r_c=2.5\,\sigma$)--- el término $\tfrac12 U_c$ es nulo en el instante de
medición y no contamina ningún canal reportado.

\subsection{Desarrollo explícito de los cuatro canales}
\label{sec:desarrollo_canales}

Con el mapeo y las dos notas anteriores ya establecidos, cada canal
$\alpha_i$ se puede desarrollar hasta quedar escrito enteramente en
términos de las ocho energías originales de la Sección~\ref{sec:2012}.

\paragraph{$\alpha_{\rm cm}$: de la cinemática de caída libre al COR.}
La bola cae desde el reposo una altura $h_0^{\rm ef}$ bajo el campo
constante $g$; toda la energía del problema nace como potencial
gravitatoria de la bola, $E_0=Mg\,h_0^{\rm ef}$. Considerando solo el
movimiento de su centro de masa (sin tocar aún los canales internos), la
velocidad de impacto satisface
\[
\tfrac12 M v_0^2 = Mg\,h_0^{\rm ef} \quad\Longrightarrow\quad
v_0=\sqrt{2g\,h_0^{\rm ef}}.
\]
Tras el primer rebote, el CM se aleja con velocidad $v_1$ y por la misma
conservación de energía del canal traslacional alcanza una altura de
apogeo $h_1$ tal que $v_1=\sqrt{2g\,h_1}$. El coeficiente de restitución
del primer rebote se define cinemáticamente como el cociente de
rapideces, $e_1\equiv v_1/v_0$ (una cantidad puramente dinámica, sin
mención de energía), de donde
\begin{equation}
e_1^2=\frac{v_1^2}{v_0^2}=\frac{2g\,h_1}{2g\,h_0^{\rm ef}}
=\frac{h_1}{h_0^{\rm ef}}.
\label{eq:e1_derivacion}
\end{equation}
Y como la energía que queda en el canal del CM tras el rebote es de nuevo
puramente potencial en el apogeo, $Mg\,h_1$, su fracción de $E_0$ es
$Mg\,h_1/(Mg\,h_0^{\rm ef})=h_1/h_0^{\rm ef}$: exactamente el mismo
cociente. Por eso ec.~(\ref{eq:acm}) escribe dos expresiones separadas por
un signo igual --- $h_1/h_0^{\rm ef}$ (definición energética) y $e_1^2$
(definición cinemática, ec.~\ref{eq:e1_derivacion}) --- que coinciden
exactamente \emph{porque} la caída y el ascenso del CM son, cada uno, un
tiro vertical simple bajo campo uniforme, sin pérdidas dentro del propio
canal del CM.

\paragraph{$\alpha_{T_B}$: por qué se evalúa justo en el apogeo.}
Desarrollando la suma de la ec.~(\ref{eq:atb}) con el tiempo explícito,
$t_1$ el instante del primer apogeo ($V_{{\rm CM},z}^B(t_1)=0$ tras la
separación):
\begin{equation}
T_B^{(1)} = \sum_{i\in\mathcal B}\tfrac12 m_i
\left\|\vec v_i(t_1)-\vec V_{\rm CM}^B(t_1)\right\|^2.
\label{eq:TB1_exp}
\end{equation}
El apogeo no es un instante arbitrario: es precisamente cuando
$K_B=\tfrac12M_B\|\vec V_{\rm CM}^B\|^2=0$ (ec.~\ref{eq:2012_KB}), de modo
que \emph{toda} la energía cinética que le queda a la bola en ese instante
es, por construcción, interna --- no hay velocidad de CM que restar ni
mezcla posible entre el canal traslacional y el vibracional. Evaluar
$\alpha_{T_B}$ en cualquier otro instante mezclaría ambos canales y
volvería ambigua la separación de la ec.~(\ref{eq:cierre}).

\paragraph{$\alpha_{\rm surf}$: energía total ganada por la superficie.}
Expandiendo la diferencia de la ec.~(\ref{eq:asurf}) en sus dos instantes,
\begin{equation}
\Delta(K_S+V_{\rm surf}) =
\big[K_S^{(1)}+V_{\rm surf}^{(1)}\big] -
\big[K_S^{(0)}+V_{\rm surf}^{(0)}\big].
\label{eq:asurf_exp}
\end{equation}
En $t=0$ la superficie fue equilibrada a $T=0$ y sus velocidades se fijan
en cero antes de soltar la bola (\texttt{velocity all set 0.0 0.0 0.0} en
\texttt{in.bounce\_cluster}), así que $K_S^{(0)}=0$ exactamente; además la
bola parte lejos de la superficie, de modo que no hay pares
bola--superficie dentro de $r_c$ y $V_{\rm surf}^{(0)}=V_S^{{\rm
paper},(0)}$ sin ninguna contaminación de $U_c$. En $t=t_1$ la bola ya se
separó y $U_c^{(1)}\approx0$ (Nota~(b) arriba), así que
$V_{\rm surf}^{(1)}=V_S^{{\rm paper},(1)}$ y (por las capas fijas de la
losa, Nota~(a)) $K_S^{(1)}=T_S^{{\rm paper},(1)}$. La ec.~(\ref{eq:asurf_exp})
se reduce entonces, en ambos extremos, a
\begin{equation}
\alpha_{\rm surf}\approx\frac{T_S^{{\rm paper},(1)}
+V_S^{{\rm paper},(1)}-V_S^{{\rm paper},(0)}}{E_0}:
\label{eq:asurf_final}
\end{equation}
la energía vibracional (cinética \emph{más} potencial) que la red de la
superficie ganó entre el instante de soltar la bola y el primer apogeo
---literalmente los fonones creados por el impacto, sin ninguna
ambigüedad de reparto de $U_c$ en ninguno de los dos extremos de la
diferencia.

\paragraph{$\alpha_{\Delta V_B}$: por qué se resta el valor inicial.}
Usando la forma explícita de $V_B$, ec.~(\ref{eq:2012_VB_exp}), en los
mismos dos instantes,
\begin{equation}
V_B^{(1)}-V_B^{(0)} =
\sum_{\substack{i<j\\ i,j\in\mathcal B}}
4\eps\left[\left(\frac{\sigma}{r_{ij}(t_1)}\right)^{12}
-\left(\frac{\sigma}{r_{ij}(t_1)}\right)^{6}\right]
-\sum_{\substack{i<j\\ i,j\in\mathcal B}}
4\eps\left[\left(\frac{\sigma}{r_{ij}(0)}\right)^{12}
-\left(\frac{\sigma}{r_{ij}(0)}\right)^{6}\right].
\label{eq:avb_exp}
\end{equation}
$V_B^{(0)}$ no es cero: es la energía de cohesión del cúmulo frío
premínimizado, con todos los pares cerca de su separación de equilibrio
$r_m$ (Sección ``Un sólido que vibra: equipartición + virial juntos'' de
\texttt{equiparticion\_y\_virial.tex}), y crece en magnitud con $N$. Si
reportáramos $V_B^{(1)}$ solo, esa cohesión ---que nada tiene que ver con
el impacto--- dominaría por completo la cifra. Restar $V_B^{(0)}$ aísla
exactamente el calentamiento potencial producido por la colisión: los
pares se alejan de $r_m$ al vibrar con más amplitud, y como el pozo de
Lennard-Jones es asimétrico (Sección ``El exceso anarmónico'' de la misma
nota), ese desplazamiento promedio es sistemáticamente positivo.

% =============================================================================
\section{Diferencias sustantivas}
\label{sec:diferencias}
% =============================================================================

\begin{enumerate}[leftmargin=1.4em]

\item \textbf{Instantánea vs.\ serie temporal.} El capítulo original
grafica las ocho energías en función continua del tiempo (su Fig.~9,
en unidades de $\eps$) para describir el mecanismo de intercambio durante
la colisión. Este estudio evalúa los canales en un único instante ---el
primer apogeo--- y los normaliza por $E_0$, sacrificando la dinámica
temporal a cambio de poder comparar de forma sistemática 210 corridas
distintas de $(N,h_0)$ mediante un solo número por canal.

\item \textbf{Un canal de superficie explícito.} El capítulo no define una
fracción de energía análoga a $\alpha_{\rm surf}$; sus energías de
superficie ($K_S$, $T_S$, $V_S$) se leen solo del gráfico temporal,
comentando que ``no contribuyen significativamente a la transferencia de
energía'' para el caso que analizan en detalle ($N\approx79$,
$h_0=11.03\,\sigma$, régimen casi elástico). Este estudio muestra que esa
conclusión \emph{no} se sostiene en general: $\alpha_{\rm surf}$ crece con
$v_{\rm imp}$ y llega a ser comparable a los otros canales a alta
velocidad (Sección ``¿A dónde se va la energía? Los cuatro canales'' de
\texttt{borrador\_rbef\_v2.tex}), lo que motivó agregar el canal como
fracción de pleno derecho, con su propio control de tamaño finito de la
losa.

\item \textbf{Cúmulos de tamaño exacto.} El capítulo construye la bola
recortando una esfera de la red (\texttt{create\_atoms region sphere}),
lo que produce un número de átomos que depende de dónde cae el centro de
la esfera respecto de la red ($N\approx79$--84 para su radio). Este
estudio corrige esto construyendo la bola por capas esféricas FCC
completas y minimizándola en vacío antes de soltarla, garantizando el
mismo $N$ exacto en todas las corridas.

\item \textbf{$\alpha_{\rm cm}$ como fracción formal.} El capítulo
reporta las alturas de rebote (su Fig.~7) pero no las convierte
explícitamente en una fracción de energía normalizada; este estudio
formaliza $\alpha_{\rm cm}=e_1^2=h_1/h_0^{\rm ef}$ como el primero de los
cuatro canales, permitiendo sumarlo directamente con los otros tres en la
ec.~(\ref{eq:cierre}).

\item \textbf{Promedio sobre orientaciones.} El capítulo reporta una
simulación por altura $h_0$ (sin barras de error). Este estudio promedia
sobre 5 orientaciones iniciales del cúmulo para cada $(N,h_0)$, reportando
media $\pm\sigma$ ---una respuesta directa a que el capítulo original no
cuantifica la variabilidad intrínseca del sistema.

\end{enumerate}

% =============================================================================
\section{Conclusión}
% =============================================================================

Las cuatro fracciones $\alpha_i$ de este estudio no son un esquema nuevo:
son las mismas ocho energías de Gutiérrez \emph{et al.} (2012), evaluadas
en un instante fijo (el primer apogeo) y renormalizadas por $E_0$ para que
sumen $1$. Las dos correcciones de bajo nivel (el reparto de $U_c$ entre
$V_B$ y $V_S$, y la definición de $K_S$) no cambian ningún resultado
reportado porque ambas se anulan en el instante de medición; las
diferencias que sí importan son conceptuales: aislar $\alpha_{\rm surf}$
como canal de pleno derecho, construir cúmulos de $N$ exacto, y promediar
sobre orientaciones para tener barras de error.

% =============================================================================
\begin{thebibliography}{9}
\bibitem{Gutierrez2012}
G. Gutiérrez, S. Davis, C. Loyola, J. Peralta, F. González, Y. Navarrete
y F. González-Wasaff, \textit{Inelastic Collisions and Hypervelocity
Impacts at Nanoscopic Level: A Molecular Dynamics Study}, cap.~12 en
\textit{Molecular Dynamics -- Theoretical Developments and Applications
in Nanotechnology and Energy}, editado por L. Wang (InTech, Rijeka, 2012).
\end{thebibliography}

\end{document}
